\documentclass{beamer}
\usepackage[utf8]{inputenc}

\usetheme{Madrid}
\usecolortheme{default}
\usepackage{amsmath,amssymb,amsfonts,amsthm}
\usepackage{txfonts}
\usepackage{tkz-euclide}
\usepackage{listings}
\usepackage{adjustbox}
\usepackage{array}
\usepackage{tabularx}
\usepackage{gvv}
\usepackage{multicol}   % <-- Add this

\usepackage{lmodern}
\usepackage{circuitikz}
\usepackage{tikz}
\usepackage{graphicx}
\usepackage{mathtools}
\setbeamertemplate{page number in head/foot}[totalframenumber]

\usepackage{tcolorbox}
\tcbuselibrary{minted,breakable,xparse,skins}



\definecolor{bg}{gray}{0.95}
\DeclareTCBListing{mintedbox}{O{}m!O{}}{%
  breakable=true,
  listing engine=minted,
  listing only,
  minted language=#2,
  minted style=default,
  minted options={%
    linenos,
    gobble=0,
    breaklines=true,
    breakafter=,,
    fontsize=\small,
    numbersep=8pt,
    #1},
  boxsep=0pt,
  left skip=0pt,
  right skip=0pt,
  left=25pt,
  right=0pt,
  top=3pt,
  bottom=3pt,
  arc=5pt,
  leftrule=0pt,
  rightrule=0pt,
  bottomrule=2pt,
  toprule=2pt,
  colback=bg,
  colframe=orange!70,
  enhanced,
  overlay={%
    \begin{tcbclipinterior}
    \fill[orange!20!white] (frame.south west) rectangle ([xshift=20pt]frame.north west);
    \end{tcbclipinterior}},
  #3,
}
\lstset{
    language=C,
    basicstyle=\ttfamily\small,
    keywordstyle=\color{blue},
    stringstyle=\color{orange},
    commentstyle=\color{green!60!black},
    numbers=left,
    numberstyle=\tiny\color{gray},
    breaklines=true,
    showstringspaces=false,
}


\title 
{5.12.5}
\date{October 3, 2025}


\author 
{Dhanush Kumar A - AI25BTECH11010}



\begin{document}
\frame{\titlepage}
\begin{frame}{Question}
Solve for $x \in \mathbb{N}$ using matrices:
$
\det\myvec{x+3 & -2 \\ -3x & 2x} = 8
$


\end{frame}
\begin{frame}{Solution}
\begin{align}
\det\myvec{x+3 & -2 \\ -3x & 2x} 
&= (x+3)(2x) - (-2)(-3x) \\
&= 2x^2 + 6x - 6x \\
&= 2x^2 \\
2x^2 &= 8 \\
x^2 &= 4 \\
x &= \pm 2 \\
x &\in \mathbb{N} \implies x = 2
\end{align}


\end{frame}



\begin{frame}[fragile]                            
\frametitle{Python code - To solve  }
	\begin{lstlisting}
import numpy as np
from sympy import symbols, Eq, solve

# Define the variable
x = symbols('x')

# Define the matrix
A = np.array([[x + 3, -2],
              [-3*x, 2*x]])

# Compute the determinant using sympy
det = (x + 3)*(2*x) - (-2)*(-3*x)

# Equation: determinant = 8
equation = Eq(det, 8)
\end{lstlisting}
\end{frame}


\begin{frame}[fragile]                            
\frametitle{Python code - To solve  }
	\begin{lstlisting}

# Solve the equation
solutions = solve(equation, x)

# Filter for natural numbers
natural_solutions = [sol for sol in solutions if sol.is_real and sol > 0]

print("All solutions:", solutions)
print("Natural number solution(s):", natural_solutions)


\end{lstlisting}
\end{frame}


\begin{frame}[fragile]                            
\frametitle{Python code - Output }                
\begin{lstlisting}
All solutions: [-2, 2]
Natural number solution(s): [2]


\end{lstlisting}

\end{frame}
\begin{frame}[fragile]                            
\frametitle{C code - Solving}
	\begin{lstlisting}
#include <stdio.h>
#include <stdlib.h>
#include <math.h>
#include "/home/dhanush-kumar-a/ee1030-2025/ai25btech11010/matgeo/5.12.5/codes/libs/matfun.h"


// Function to solve determinant equation det([[x+3, -2], [-3*x, 2]]) = 8
double solve_det_equation() {
    double **A = createMat(2, 2);

    // We'll solve the equation 2*x^2 = 8, but let's use matrix det
    // Represent matrix as [[x+3, -2], [-3*x, 2]]
    // We'll solve symbolically by using quadratic formula
    double a = 2.0;   // coefficient of x^2
    double b = 0.0;   // coefficient of x
    double c = -8.0;  // constant term
    double **roots = Matquad(a, b, c); // Solve 2*x^2 - 8 = 0
\end{lstlisting}

\end{frame}
\begin{frame}[fragile]                            
\frametitle{C code - Solving}
	\begin{lstlisting}

    double x = roots[0][0]; // Return positive root (x ∈ N)
    
    // Free allocated memory
    free(A);
    free(roots);

    return x;
}

// Entry point for Python ctypes
double solve_det_ctypes() {
    return solve_det_equation();
}



\end{lstlisting}
\end{frame}

	
\begin{frame}[fragile]                              
	\frametitle{Python code -Solving using c Funtion} 
	\begin{lstlisting}
import ctypes

lib = ctypes.CDLL('./libsolve_det.so')  # Shared library compiled from above C code
lib.solve_det_ctypes.restype = ctypes.c_double

x = lib.solve_det_ctypes()
print("Solution x ∈ N:", x)


\end{lstlisting}
\end{frame}

	
\begin{frame}[fragile]                              
	\frametitle{Python code-Output} 
	\begin{lstlisting}
Solution x ∈ N: 2.0

\end{lstlisting}                               
\end{frame}


	


\end{document}




